\chapter{Common Postprocessing Tasks}

\section{Data Exchange with CFS++}

\subsection{ANSYS Classic}
\subsection{GiD}
\subsection{Gmsh}
\subsection{ANSYS CFDPost}
\subsection{GMV}

\subsection{EnSight}

EnSight  supports the  ANSYS .rst  file format  and can  therefore  read files
generated   by  our   ANSYS  output   writer  (cf.   Figs.  \ref{fig:ensight},
\ref{fig:cfs_soft_landscape}).  In order for  EnSight to recognize our regions
as parts, it is  necessary to specify the fact that our  regions get mapped to
material  ids in the  .rst files  to its  ANSYS reader.  One has  therefore to
switch to  the advanced  interface and  then select 'Use  Material Id'  in the
'Regular          Part          Creation         Convention'          combobox
(cf. Fig. \ref{fig:ensight_rst_mat_to_part}).

\begin{figure}[htb] 
  \centering 
  \subfloat[The EnSight main window.]{
    \label{fig:ensight}
    \includegraphics[height=7.5cm]{pics/ensight_man} }
  \qquad
  \subfloat[Let EnSight make parts out of materials.]{
    \label{fig:ensight_rst_mat_to_part}
    \includegraphics[height=7.56cm]{pics/ensight_rst_mat_to_part}
  } 
  \caption{Postprocessing with EnSight.}
\end{figure}

\subsection{Tecplot}
\subsection{VisIt}
\subsection{ParaView}
\subsection{Matlab}
\subsection{Gnuplot}

\section{Creating Color Mapped Plots}

\subsection{ParaView}
\url{https://visualization.hpc.mil/wiki/Coloring_Data_in_Paraview}
\url{http://paraview.org/OnlineHelpCurrent/Cut.html}

\subsection{GiD}
\url{http://gid.cimne.upc.es/component/manual/24170?id=24170}

\url{https://lse17.e-technik.uni-erlangen.de:2001/docu/tutorial/html/cfsManual-html.html}

\section{Creating Slices}

\subsection{Tecplot}
\url{http://www.tecplot.com/Portals/0/Tutorials/Slice.htm}

\subsection{ParaView}
\url{https://visualization.hpc.mil/wiki/Paraview_Slice}

\subsection{GiD}
\url{http://gid.cimne.upc.es/component/manual/24170?id=24170}

\subsection{EnSight}
\url{https://visualization.hpc.mil/wiki/Working_with_Slicing_Planes}

\section{Creating Isocontuour and Isosurface Plots}

\subsection{ParaView}
\url{https://visualization.hpc.mil/wiki/Paraview_Contours_and_Isosurfaces}

\url{http://paraview.org/OnlineHelpCurrent/Contour.html}

\subsection{GiD}
\url{http://gid.cimne.upc.es/component/manual/24164?id=24164}
\url{https://lse17.e-technik.uni-erlangen.de:2001/docu/tutorial/html/cfsManual-html.html}


\subsection{Tecplot}
\url{http://www.tecplot.com/Portals/0/Tutorials/Contour.htm}

\url{http://www.tecplot.com/Portals/0/Tutorials/3DContour.htm}

\url{http://www.tecplot.com/Portals/0/Tutorials/Isosurface.htm}

\subsection{EnSight}
\url{https://visualization.hpc.mil/wiki/Working_with_Isosurfaces}

\section{Visualizing Deformed Geometry}

\subsection{ParaView}
\url{http://paraview.org/OnlineHelpCurrent/WarpVector.html}
\url{http://paraview.org/OnlineHelpCurrent/WarpScalar.html}

\subsection{GiD}
\url{https://lse17.e-technik.uni-erlangen.de:2001/docu/tutorial/html/cfsManual-html.html}

\subsection{GMV}

Although GMV is  very good in visualizing scalar and  vector data with contour
plots, isosurfaces  and glyphs,  it has unfortunately  no built-in  ability to
visualize  deformed geometries.   This  situation will  not  change in  future
versions  of GMV,  since at  LANL there  is no  need for  visualizing deformed
meshes.

The only  possibilty is to  re-write gmv files with  mechanical displacements,
where  the  displacements are  added  to  the  nodal coordinates  as  offsets.
Fortunately, there exists  a GMV writer/reader-library\footnote{GMV I/O source
codes:  \url{http://www.generalmeshviewer.com/io_util.html}},  which was  used
for     the    program     gmv-addDisp\footnote{GMV     displacement    tools:
\url{https://lse17.e-technik.uni-erlangen.de:2001/files/precompiled/GMV/gmv-tools.tgz}}.
It can be used  to read in a GMV file, add the  displacements as offset to the
nodal coordinates and write a new result file with a new name.  The program is
also located in /home/data/programs/bin.

\noindent The syntax is as follows:
\begin{verbatim}
  gmv-addDisp [in-filename] [out-filename] [relative/absolute] [factor]
\end{verbatim}

\noindent with the following parameters

\begin{itemize}
    \item in-filename: Name of the input gmv file
    \item out-filename: Name of the resulting deformed gmv file
    \item relative/absolute:  Indicates if  relative (= 1)  or absolute  (= 0)
scaling  should be  used.  For relative  scaling,  factor is  relative to  the
bounding box of the geometry
    \item factor: Scaling factor for the nodal displacements
\end{itemize}

A example run, where a given file  bar.gmv is deformed by a relative factor of
10\% and written to the new file def-bar.gmv, would look like this:

\begin{verbatim}
user@lse90: gmv-addDisp bar.gmv def-bar.gmv 1 0.1

Determine bounding box.
Maximum displacement: 5.568630e-08
Abs. Factor: 2.514083e+07
Converting nodev.
Converting cells.
Converting materials.
Converting groups.
Converting variables.
Converting problem time.
Converting cycle number.
Success!
user@lse90:
\end{verbatim}

\begin{figure}[htb] 
  \centering 
  \subfloat[Undeformed bar.]{
    \label{fig:gmv_def_bar}
    \includegraphics[width=6cm]{Bar-undef} }
  \qquad
  \subfloat[Deformed bar.]{
    \label{fig:gmv_undef_bar}
    \includegraphics[width=6cm]{Bar-def}
  } 
  \caption{The  maximum deformation  of  the bar  can  be also  seen from  the
output.}
  \label{fig:gmv_bars}
\end{figure}

For  a larger  number of  files the  script gmv-deform.sh  can be  used, which
iterates over a given list of GMV files and creates deformed output files with
the prefix def-. Therefore an absolute factor has to be provided:

\begin{verbatim}
  gmv-deform.sh [absolute factor] '[list of GMV files]'
\end{verbatim}

NOTE: Don't  forget to put the  ' ' around  the filenames in order  to prevent
filename-expansion by the shell!


\section{Visualizing Vector and Tensor Data with Glyphs}

\subsection{ParaView}
\url{https://visualization.hpc.mil/wiki/Glyph}

\subsection{GiD}
\url{https://lse17.e-technik.uni-erlangen.de:2001/docu/tutorial/html/cfsManual-html.html}

\subsection{Tecplot}
\url{http://www.tecplot.com/Portals/0/Tutorials/Vector.htm}

\section{Selecting Elements or Nodes}

\section{Restoring Visualizations from State Files}


\section{Creating Animations}

\subsection{GiD}
\url{https://lse17.e-technik.uni-erlangen.de:2001/docu/tutorial/html/cfsManual-html.html}


\section{Obtaining Signal Data from Mesh Postprocessors}

\subsection{GiD}
\url{https://lse17.e-technik.uni-erlangen.de:2001/docu/tutorial/html/cfsManual-html.html}
\url{http://gid.cimne.upc.es/component/manual/24180?id=24180}

\section{Plotting Output/History Signals}

\section{Creating Directivity Plots}

\subsection{Matlab/Octave}

\subsection{Gnuplot}